% DynamicScaler Implemented Changes Report
% Based on comparison: https://github.com/sh-Lin/DynamicScaler/compare/main...AnshRaj112:DynamicScaler:main

\documentclass[11pt,a4paper]{article}

\usepackage[utf8]{inputenc}
\usepackage[T1]{fontenc}
\usepackage{lmodern}
\usepackage[margin=1in]{geometry}
\usepackage{amsmath,amssymb,amsfonts}
\usepackage{hyperref}
\usepackage{xcolor}
\usepackage{listings}
\usepackage{booktabs}
\usepackage{enumitem}
\usepackage{float}
\usepackage{titlesec}

\hypersetup{
    colorlinks=true,
    linkcolor=blue!60!black,
    citecolor=green!40!black,
    urlcolor=cyan!50!black,
}

\lstset{
    language=Python,
    basicstyle=\ttfamily\footnotesize,
    keywordstyle=\color{blue},
    commentstyle=\color{gray},
    stringstyle=\color{red!70!black},
    breaklines=true,
    frame=single,
    numbers=left,
    numberstyle=\tiny\color{gray},
    showstringspaces=false,
    tabsize=4,
}

\definecolor{enhance}{RGB}{0,100,0}
\newcommand{\enhance}[1]{\textcolor{enhance}{\textbf{#1}}}
\newcommand{\file}[1]{\texttt{#1}}
\newcommand{\fn}[1]{\texttt{#1}}

\title{
    \textbf{DynamicScaler: Implemented Code, Math, and Low-Memory Improvements} \\[4pt]
    \large From Original CVPR 2025 Release to Optimized Fork
}
\author{Engineering Change Log and Rationale}
\date{\today}

\begin{document}
\maketitle

\begin{abstract}
This document summarizes the concrete changes introduced in the fork of \textbf{DynamicScaler} relative to the original main branch, as captured by the comparison
\href{https://github.com/sh-Lin/DynamicScaler/compare/main...AnshRaj112:DynamicScaler:main}{main\,$\dots$\,main}.
For each modified component we describe:
(i) the \emph{previous behavior} in code and mathematics,
(ii) the \emph{new implementation},
and (iii) the \emph{purpose and expected effect} on image quality, temporal stability, geometry, and low-memory operation (e.g.\ Tesla P100 with 16\,GB VRAM).
\end{abstract}

\tableofcontents
\newpage

%==============================================================================
\section{High-Level Overview of Changes}
%==============================================================================

The recent commits introduce three main classes of improvements:
\begin{enumerate}
    \item \textbf{Geometric and sampling correctness} in the panorama tensor utilities (sphere and ring), improving projection accuracy and multi-view blending.
    \item \textbf{Diffusion pipeline behavior} changes, including DDIM-aware re-noising, temporal guidance in the sphere pipeline, smoother overlap/merge schedules, and post-decode temporal smoothing.
    \item \textbf{Low-memory and efficiency features}, such as adaptive tiling, caching of invariant computations, resolution and step reductions in ``low-memory'' mode, and an explicit 2$\times$-upscale safety switch.
\end{enumerate}

The analysis documents \file{enhancement\_analysis.tex} and \file{comprehensive\_improvements.tex} previously proposed many of these changes.
The current fork \emph{implements} a large subset of those proposals; this report focuses only on what has actually changed in code.

%==============================================================================
\section{Changes in \file{gen\_pano\_360.py}}
%==============================================================================

\subsection{New Configuration Flags}

\paragraph{Previous.}
The script did not expose:
\begin{itemize}[nosep]
    \item A temporal guidance scale for the sphere pipeline.
    \item A global ``low-memory'' switch.
    \item A hard safety guard to skip the 2$\times$ upscale stage on 16\,GB GPUs.
\end{itemize}

\paragraph{New.}
The argument dataclass \fn{VArgs} now has:
\begin{lstlisting}
temporal_guidance_scale: Optional[float] = None
low_memory: bool = False
\end{lstlisting}
and there is a module-level constant:
\begin{lstlisting}
SKIP_2X_UPSCALE_TO_AVOID_OOM = True
\end{lstlisting}
The value of \fn{temporal_guidance_scale} is threaded into the sphere panorama pipeline, and \fn{low\_memory} drives multiple conservative choices tailored to GPUs such as Tesla P100 (16\,GB).

\paragraph{Purpose.}
\begin{itemize}
    \item \textbf{Temporal guidance scale} enables classifier-free \emph{temporal} guidance in the sphere pipeline, improving motion smoothness and frame consistency.
    \item \textbf{Low-memory} mode centralizes all P100--class heuristics so that a single flag can be used to avoid out-of-memory errors.
    \item \textbf{2$\times$-upscale guard} prevents the most memory-intensive VAE decode from running by default on 16\,GB cards.
\end{itemize}

\subsection{Low-Memory Denoising Resolution and Steps}

\paragraph{Previous.}
The base video diffusion sampler always ran with:
\[
    \text{height}=512,\quad \text{width}=1024,\quad \text{num\_steps}=48.
\]

\paragraph{New.}
At the script entry point, these are now conditioned on \fn{vargs.low\_memory}:
\begin{itemize}
    \item If \fn{low\_memory=False}:
    \[
        \text{height}=512,\quad \text{width}=1024,\quad \text{num\_steps}=48.
    \]
    \item If \fn{low\_memory=True}:
    \[
        \text{height}=256,\quad \text{width}=384,\quad \text{num\_steps}=24.
    \]
\end{itemize}

\paragraph{Purpose.}
This reduces both spatial and temporal compute requirements for the UNet 3D:
the memory footprint scales roughly with \(H \times W \times T\).
Halving each spatial dimension and halving the number of steps makes the sphere stage tractable on 16\,GB VRAM, especially when combined with later tiling and overlap adjustments.

\subsection{Low-Memory Equirectangular Base Resolution}

\paragraph{Previous.}
The base equirectangular panorama resolution was fixed to
\[
    W_\text{eq}=1024,\quad H_\text{eq}=512,
\]
and the sphere stage worked on the doubled resolution (due to internal scaling) regardless of GPU.

\paragraph{New.}
Now:
\begin{itemize}
    \item If \fn{low\_memory=False}, the same \(1024\times512\) baseline is preserved.
    \item If \fn{low\_memory=True}, the sphere stage instead uses
    \[
        W_\text{eq}=512,\quad H_\text{eq}=256.
    \]
\end{itemize}

\paragraph{Purpose.}
This directly lowers the size of the 5D latent tensor \((B,C,T,H,W)\) in the sphere stage by a factor of \(4\times\).
Although this reduces raw resolution, the later planar/1$\times$ and optional 2$\times$ pathways can restore apparent resolution while preserving temporal stability.

\subsection{Latitude-Adaptive Theta Sampling}

\paragraph{Previous.}
The mapping \fn{phi\_theta\_dict} used a fixed number of azimuthal samples for most latitudes:
\begin{lstlisting}
phi_theta_dict = {
    90:  [0],
    -90: [0],
    75:  [360*t//phi_num for t in range(phi_num)],
    ...
    0:   [360*t//phi_num for t in range(phi_num)],
}
\end{lstlisting}
This oversampled near the poles (where longitude lines converge) and undersampled near the equator.

\paragraph{New.}
The number of theta samples per latitude \(\phi\) is now:
\begin{equation}
    N_\theta(\phi) = \max\bigl(1,\;\lfloor N_0 \cos(\phi \cdot \pi/180)\rfloor\bigr),
\end{equation}
where \(N_0=\texttt{phi\_num}\) is the base equatorial count.
For example:
\begin{lstlisting}
def n_theta_for_phi(phi_deg, n0):
    n = max(1, int(n0 * math.cos(phi_deg * math.pi / 180.0)))
    return max(1, n)

phi_theta_dict = {
    90:  [0],
    -90: [0],
    75:  [360 * t // n_theta_for_phi(75, phi_num) for t in range(n_theta_for_phi(75, phi_num))],
    ...
    0:   [360 * t // phi_num for t in range(phi_num)],
}
\end{lstlisting}

\paragraph{Purpose.}
This approximates uniform solid-angle coverage:
at latitude \(\phi\), the circumference of a parallel is proportional to \(\cos\phi\),
so the number of theta samples should also scale with \(\cos\phi\).
As a result:
\begin{itemize}
    \item Polar regions are no longer oversampled with nearly redundant views.
    \item Equatorial regions receive more relative sampling, increasing scene richness where most content lives.
    \item Overall compute is reduced for the same perceived coverage.
\end{itemize}

\subsection{Cosine-Annealed Overlap Ratios}

\paragraph{Previous.}
Both the 1$\times$ and 2$\times$ overlap ratios used a two-point step function between \(0.75\) and \(0.5\),
implemented via integer indexing into a short list.
This produced an abrupt mid-schedule change in overlap.

\paragraph{New.}
For both 1$\times$ and 2$\times$ paths a cosine-annealed schedule is used:
\begin{equation}
    r(t) = r_{\max} - \frac{r_{\max}-r_{\min}}{2}
    \left(1 - \cos\frac{\pi t}{T-1}\right),
    \quad r_{\max}=0.75,\; r_{\min}=0.5,
\end{equation}
for \(t=0,\dots,T-1\) where \(T\) is \fn{num\_inference\_steps}.

\paragraph{Purpose.}
The overlap between shift windows changes smoothly across timesteps, avoiding a sudden change in effective receptive field.
Empirically this reduces temporal seam artifacts and flickering at the point where overlap would otherwise jump from \(0.75\) to \(0.5\).

\subsection{Cosine Decay of Merge-Previous-Denoised Ratio}

\paragraph{Previous.}
The list \fn{merge\_prev\_denoised\_ratio\_list} decayed linearly from the maximum ratio down to zero over a fixed number of merge steps.
Linear decay treats early and late steps equally, even though late steps are mostly detail refinement.

\paragraph{New.}
The merged ratio now uses a cosine decay:
\begin{equation}
    \lambda(t) = \frac{\lambda_{\max}}{2} \left(1 + \cos\frac{\pi t}{T_\text{merge}}\right),\qquad 0\le t\le T_\text{merge},
\end{equation}
and is set to zero thereafter.

\paragraph{Purpose.}
This keeps strong blending early, when structure is coarse and past frames provide valuable guidance, then rapidly decreases blending during the fine-detail phase.
It improves sharpness and reduces over-smoothed textures near convergence, while maintaining loop consistency.

\subsection{Latent Resize Mode: Bilinear Instead of Nearest}

\paragraph{Previous.}
The downsampled sphere SW latent used nearest-neighbor interpolation:
\begin{lstlisting}
downsampled_sphere_SW_latent = resize_video_latent(
    input_latent=sphere_SW_latent.clone(),
    mode="nearest", ...)
\end{lstlisting}

\paragraph{New.}
The call now forces bilinear mode:
\begin{lstlisting}
downsampled_sphere_SW_latent = resize_video_latent(
    input_latent=sphere_SW_latent.clone(),
    mode="bilinear", ...)
\end{lstlisting}

\paragraph{Purpose.}
Nearest-neighbor upscaling/downscaling in latent space introduces blockiness and aliasing.
Using bilinear interpolation creates smoother latents which decode to cleaner images, especially when combined with the optional Gaussian blur introduced in \file{diffusion\_utils.py}.

\subsection{Emptying GPU Cache Between Stages}

\paragraph{Previous.}
No explicit cache clearing occurred between sphere and planar stages, leaving peak memory potentially dominated by residual allocations from prior steps.

\paragraph{New.}
When \fn{vargs.low\_memory=True}, calls to:
\begin{lstlisting}
torch.cuda.empty_cache()
gc.collect()
\end{lstlisting}
are inserted after key stages.

\paragraph{Purpose.}
This reduces VRAM fragmentation and releases memory from temporary tensors that are no longer needed,
making it more likely that subsequent stages (especially 1$\times$ planar denoising) will fit into the 16\,GB budget.

%==============================================================================
\section{Sphere Panorama Pipeline: \file{i2v\_sphere\_panorama\_pipeline.py}}
%==============================================================================

\subsection{Temporal Guidance in Sphere Pipeline}

\paragraph{Previous.}
The sphere panorama pipeline only used text+image classifier-free guidance.
Temporal guidance---already implemented in the DDIM sampler for planar videos---was not applied in the sphere multi-view path.

\paragraph{New.}
A scalar \fn{temporal\_guidance\_scale} is now accepted by \fn{basic\_sample\_shift\_shpere\_panorama} and,
when nonzero, the model performs two forward passes per view:
\begin{itemize}
    \item \(\epsilon_t^\text{full}\): standard prediction with temporal attention enabled.
    \item \(\epsilon_t^\text{image}\): prediction with \fn{no\_temporal\_attn=True}, analogous to a per-frame image model.
\end{itemize}
The temporally guided prediction is:
\begin{equation}
    \epsilon_t^\text{temp} = \epsilon_t^\text{full}
    + \lambda_\text{temp} \bigl(\epsilon_t^\text{full} - \epsilon_t^\text{image}\bigr),
\end{equation}
where \(\lambda_\text{temp} = \texttt{temporal\_guidance\_scale}\).

\paragraph{Purpose.}
This is the temporal analogue of standard CFG:
it emphasizes noise components that are consistent across frames (via the difference between full and image-only predictions).
The effect is improved motion coherence and reduced ``frame-wise drifting'' across time while preserving prompt adherence.

\subsection{Theta Offset: Float Division to Avoid Angular Gaps}

\paragraph{Previous.}
Theta offsets used integer division:
\begin{equation}
    \theta_\text{offset} = (i \bmod N)\cdot \bigl\lfloor\text{FoV}/N\bigr\rfloor,
\end{equation}
which can leave unvisited angular gaps when \(\text{FoV}\) is not divisible by \(N\).

\paragraph{New.}
The offset now uses floating-point division:
\begin{equation}
    \theta_\text{offset} = (i \bmod N)\cdot \frac{\text{FoV}}{N}.
\end{equation}

\paragraph{Purpose.}
This ensures that the sequence of windows covers the intended FoV evenly without holes,
improving dynamic degree and reducing un-sampled stripes in the equirectangular panorama.

\subsection{Caching: VAE Encoding and Text Embeddings}

\paragraph{Previous.}
\begin{itemize}
    \item With \fn{paste\_on\_static=True}, the same static panorama image was re-encoded by the VAE at \emph{every timestep}, even though it never changed.
    \item Elevation-specific text prompts from \fn{phi\_prompt\_dict} were re-embedded via CLIP at every view and timestep.
\end{itemize}

\paragraph{New.}
\begin{enumerate}
    \item The VAE encoding of the static panorama is computed \emph{once} before the denoising loop and reused:
    \[
        \text{frame\_0\_latent} \leftarrow \text{cached encode}.
    \]
    \item A dictionary \(\phi\mapsto \text{text\_embedding}\) is precomputed before the loop, and the cached embeddings are reused for all timesteps and theta positions at that elevation.
\end{enumerate}

\paragraph{Purpose.}
Both changes dramatically reduce redundant compute:
\begin{itemize}
    \item The tiled VAE encode of a large panorama is expensive; eliminating \(\mathcal{O}(T)\) identical encodes saves seconds to minutes per sample.
    \item CLIP text encoders are also heavy; caching per elevation slashes the number of embeddings by roughly \(\mathcal{O}(T \cdot N_\theta(\phi))\).
\end{itemize}
This is particularly important in low-memory and low-throughput environments.

\subsection{Robust Shift-Window Parameterization in Low-Resolution Settings}

\paragraph{Previous.}
The shift-window code computed overlap ratios and step sizes analytically and enforced:
\[
    0 \leq \text{overlap\_ratio} < 1, \qquad \text{latent\_offset\_step\_size} \geq 1
\]
with hard assertions.
On very small resolutions (as used in low-memory mode), rounding could yield ratios slightly outside valid ranges or step sizes below 1, triggering assertion failures.

\paragraph{New.}
The code now:
\begin{itemize}
    \item Clamps overlap ratios into \([0, 0.9999]\) instead of asserting.
    \item Ensures the latent offset step size is at least 1 when there is more than one window.
\end{itemize}

\paragraph{Purpose.}
This makes the sampler robust to aggressive downscaling and low-memory presets.
Instead of crashing, the pipeline will adaptively clamp to the nearest valid configuration, preserving the overall shift-window pattern.

\subsection{Adaptive Tiled VAE Encoding with Cosine-Tapered Blending}

\paragraph{Previous.}
The tiled VAE encoder:
\begin{itemize}
    \item Used a fixed \(4\times 4\) tile grid with fixed overlaps.
    \item Uniformly averaged overlapping latent values (equally weighted tiles),
    making seams visible when the encoder's receptive field extended beyond the overlap.
\end{itemize}

\paragraph{New.}
The encoder:
\begin{enumerate}
    \item Chooses tile counts \((h_\text{tiles}, w_\text{tiles})\) based on resolution (aiming for \(\approx64\) latent pixels per tile) and derives overlaps adaptively.
    \item Applies a separable cosine-tapered weight in the overlap regions:
    \begin{equation}
        W(x,y) = w_h(x)\,w_w(y),
    \end{equation}
    where \(w_h\) and \(w_w\) decay smoothly from \(1\) inside tiles toward the overlap borders using:
    \[
        w(d) = \tfrac{1}{2}\bigl(1+\cos(\pi d / O)\bigr), \quad d\in[0,O],
    \]
    with \(O\) the overlap in latent pixels.
    \item The global latent is then:
    \[
        L_{\text{fused}}(x,y) =
        \frac{\sum_k W_k(x,y)L_k(x,y)}{\sum_k W_k(x,y)}.
    \]
\end{enumerate}

\paragraph{Purpose.}
This accomplishes two goals:
\begin{itemize}
    \item \textbf{Memory scaling}: more, smaller tiles reduce per-encode memory and allow high-resolution images to be encoded on 16\,GB GPUs.
    \item \textbf{Visual quality}: cosine weighting suppresses hard seams between tiles by softly fading their contributions, which is especially important for smooth skies and large surfaces typical in panoramic content.
\end{itemize}

%==============================================================================
\section{Scheduler Re-noising: \file{pipeline/scheduler.py}}
%==============================================================================

\subsection{From Direct Noise Injection to DDIM-Aware Re-noising}

\paragraph{Previous formulation.}
The method \fn{re\_noise(x\_a, step\_a, step\_b)} implemented:
\begin{equation}
    x_b = \sqrt{\frac{\bar{\alpha}_b}{\bar{\alpha}_a}}\,x_a
          + \sqrt{1 - \frac{\bar{\alpha}_b}{\bar{\alpha}_a}}\,\tilde{\epsilon},
\end{equation}
with fresh Gaussian noise \(\tilde{\epsilon}\sim\mathcal{N}(0,I)\).
This ``direct'' re-noising does not preserve the internal DDIM noise trajectory.

\paragraph{New formulation.}
The method now optionally accepts \fn{noise\_pred\_a} and, if provided, performs DDIM-inversion-style re-noising:
\begin{align}
    \hat{x}_0 &= \frac{x_a - \sqrt{1-\bar{\alpha}_a}\,\hat{\epsilon}_a}{\sqrt{\bar{\alpha}_a}},\\
    x_b &= \sqrt{\bar{\alpha}_b}\,\hat{x}_0
           + \sqrt{1-\bar{\alpha}_b-\sigma_b^2}\,\hat{\epsilon}_a
           + \sigma_b\,\epsilon,
\end{align}
where \(\hat{\epsilon}_a=\texttt{noise\_pred\_a}\), \(\epsilon\sim\mathcal{N}(0,I)\), and \(\sigma_b\) is drawn from the DDIM sigma schedule.
If \fn{noise\_pred\_a} is not provided, the method falls back to the original direct formulation.

\paragraph{Purpose.}
The DDIM-aware path preserves structural information when changing noise levels between scales (e.g.\ between base and upscale stages).
Instead of injecting independent noise, the method reuses the estimated noise prediction, leading to more coherent refinements and better preservation of sharp details across stages.

%==============================================================================
\section{Latent Resizing and Spatial Smoothing: \file{diffusion\_utils.py}}
%==============================================================================

\subsection{Gaussian-Smoothed Latent Upsampling}

\paragraph{Previous.}
Latent resizing used \fn{F.interpolate} without any anti-aliasing or smoothing,
and there was no way to lightly blur upsampled latents before re-noising.

\paragraph{New.}
The signature of \fn{resize\_video\_latent} now includes \fn{gaussian\_sigma}, and when this parameter is positive,
the upsampled latent is processed by a spatial Gaussian filter:
\begin{equation}
    L_\text{smooth} = G_\sigma * L_\text{upsampled},
\end{equation}
where \(G_\sigma\) is a 2D Gaussian kernel applied per frame (via a grouped 2D convolution over the 5D tensor).

\paragraph{Purpose.}
This optional smoothing suppresses high-frequency artifacts introduced by interpolation,
especially before a re-noising step in the upscale path.
It yields crisper yet less noisy final outputs by slightly regularizing the latent manifold.

%==============================================================================
\section{Temporal Flicker and Video Saving: \file{loop\_merge\_utils.py}}
%==============================================================================

\subsection{Temporal Gaussian Smoothing}

\paragraph{Previous.}
Decoded video latents were converted to images frame-by-frame with no temporal filtering:
any per-frame noise or slight variation directly appeared as flicker.

\paragraph{New.}
\file{loop\_merge\_utils.py} introduces:
\begin{itemize}
    \item A 1D Gaussian smoothing function along the time dimension:
    \[
        x^\prime_t = \sum_{\tau} k(t-\tau)\,x_\tau,
    \]
    where \(k\) is a normalized Gaussian kernel in time.
    \item An extended \fn{save\_decoded\_video\_latents} that:
    \begin{enumerate}
        \item Optionally applies this temporal smoothing to the decoded frame tensors (controlled by \fn{temporal\_smooth\_sigma}).
        \item Decodes and processes frames in configurable mini-batches for improved GPU utilization.
    \end{enumerate}
\end{itemize}

\paragraph{Purpose.}
A small temporal Gaussian (\(\sigma\approx 0.5\)–\(1.0\)) reduces inter-frame flickering and stabilizes subtle textures (e.g.\ grass, water, clouds) without significantly smearing motion.
Batch processing is an efficiency improvement that does not change output semantics.

%==============================================================================
\section{Panorama Geometry and Blending: \file{panorama\_tensor\_utils.py} and \file{ring\_panorama\_tensor\_utils.py}}
%==============================================================================

\subsection{Equirectangular Longitude Normalization}

\paragraph{Previous.}
Longitude-to-pixel mapping used:
\begin{equation}
    u = \frac{\text{lon}}{2\pi}(W-1),
\end{equation}
so the rightmost pixel column mapped to \(\text{lon}=2\pi\), which is geometrically identical to \(\text{lon}=0\), duplicating one column.

\paragraph{New.}
Now:
\begin{equation}
    u = \frac{\text{lon}}{2\pi}W,
\end{equation}
with wrap-around handled by modular arithmetic downstream.

\paragraph{Purpose.}
This spreads \(W\) pixels uniformly over \([0,2\pi)\) without duplication, improving use of horizontal resolution and reducing seam artifacts at the 360° wrap.

\subsection{Centered Pixel Grid for Gnomonic Projection}

\paragraph{Previous.}
The image grid for view extraction used:
\[
    x = -\frac{W}{2},\dots,\frac{W}{2}-1,\qquad
    y = -\frac{H}{2},\dots,\frac{H}{2}-1,
\]
introducing a half-pixel bias when \(W\) or \(H\) are even.

\paragraph{New.}
The grid is now symmetric about the optical center:
\begin{equation}
    x_i = \frac{2i-(W-1)}{2},\quad
    y_j = \frac{2j-(H-1)}{2},
\end{equation}
implemented via:
\begin{lstlisting}
x = torch.linspace(-(width-1)/2, (width-1)/2, steps=width)
y = torch.linspace(-(height-1)/2, (height-1)/2, steps=height)
\end{lstlisting}

\paragraph{Purpose.}
This removes the systematic 0.5\,px offset in both axes, improving geometric consistency between extracted perspective views and the equirectangular panorama across all latitudes.

\subsection{Nearest-Neighbor Sampling: Floor vs Round}

\paragraph{Previous.}
Nearest-neighbor sampling computed indices using \(\lfloor u\rfloor, \lfloor v\rfloor\),
which is effectively a ``round down'' scheme and biases sampling toward the negative direction.

\paragraph{New.}
Indices now use standard nearest rounding:
\begin{equation}
    u_0 = \operatorname{round}(u),\qquad v_0 = \operatorname{round}(v).
\end{equation}

\paragraph{Purpose.}
This reduces average reprojection error (from \(\approx 0.5\) to \(\approx 0.25\) pixels under uniform sub-pixel offsets),
resulting in slightly sharper and more symmetric view reconstructions, especially important for small details at high zoom.

\subsection{Bilinear Splatting: Vectorization and Blending}

\paragraph{Previous.}
Bilinear write-back from a view to the panorama:
\begin{itemize}
    \item Used Python loops over batch and channel, with repeated \fn{index\_add\_} calls.
    \item Overwrote the panorama content wherever the weight sum was non-zero:
    \[
        P(u,v) \leftarrow \frac{\sum_k w_k V_k}{\sum_k w_k},
    \]
    ignoring any existing contributions.
\end{itemize}

\paragraph{New.}
The implementation:
\begin{enumerate}
    \item Vectorizes accumulation using \fn{scatter\_add\_} over the flattened panorama for all batches and channels.
    \item Introduces a blend factor \(\alpha\in[0,1]\) so that:
    \[
        P_\text{new} =
        \alpha \cdot \frac{\sum_k w_k V_k}{\sum_k w_k}
        + (1-\alpha) \cdot P_\text{old}
    \]
    at covered pixels.
\end{enumerate}

\paragraph{Purpose.}
This improves both speed and quality:
\begin{itemize}
    \item \textbf{Speed}: vectorization removes nested Python loops, reducing overhead for large view batches.
    \item \textbf{Quality}: blending with the existing panorama allows a smoother fusion of overlapping views and multi-pass contributions, especially when using \(\alpha<1\) for soft updates.
\end{itemize}

%==============================================================================
\section{Preprocessing and Image Loading: \file{precast\_latent\_utils.py} and \file{tensor\_utils.py}}
%==============================================================================

\subsection{Unified Resizing and Normalization}

\paragraph{Previous.}
Two pipelines coexisted:
\begin{enumerate}
    \item \file{precast\_latent\_utils.py} used \fn{transforms.Resize} (default interpolation) plus a no-op \fn{Normalize(mean=0, std=1)} followed by \((x*2-1)\).
    \item \file{tensor\_utils.py} used \fn{cv2.resize(..., INTER\_LINEAR)} and then normalized with \((x/255-0.5)*2\).
\end{enumerate}
The mismatch in interpolation and redundant normalization caused small but systematic differences in pixel statistics.

\paragraph{New.}
The two code paths are now aligned:
\begin{itemize}
    \item \file{precast\_latent\_utils.py} uses:
    \[
        \text{Resize(image\_size, antialias=True)},\quad
        \text{ToTensor},\quad
        (x/255 - 0.5)\cdot 2.
    \]
    \item \file{tensor\_utils.py} chooses interpolation based on scale factor:
    \begin{itemize}[nosep]
        \item \fn{INTER\_AREA} for downscaling.
        \item \fn{INTER\_LANCZOS4} for upscaling.
    \end{itemize}
    and also normalizes using \((x/255.0 - 0.5)\cdot 2.0\).
\end{itemize}

\paragraph{Purpose.}
This unifies the preprocessing semantics:
all conditioning images are now consistently mapped to \([-1,1]\) with high-quality interpolation.
That stability feeds into the VAE and diffusion model, reducing unexpected differences between input paths.

%==============================================================================
\section{Quality Metrics Stub: \file{quality\_metrics.py}}
%==============================================================================

\subsection{Q-Align and CLIP Score Hooks}

\paragraph{Previous.}
The codebase did not have explicit entry points for computing Q-Align (image/video) or CLIP-based quality metrics.

\paragraph{New.}
The new module \file{quality\_metrics.py} defines placeholder functions:
\begin{itemize}
    \item \fn{compute\_qalign\_image\_score}: for per-frame image quality scoring.
    \item \fn{compute\_qalign\_video\_score}: for video-level quality.
    \item \fn{compute\_clip\_score}: for prompt-frame alignment via CLIP cosine similarity.
\end{itemize}
These functions accept external models or APIs but currently return \fn{None} by default.

\paragraph{Purpose.}
This is a structural change preparing the pipeline for integration of learned quality metrics:
Q-Align and CLIP scores can be used for evaluation, ablations (e.g.\ comparing overlap schedules), or best-of-$k$ selection without modifying the core sampler logic.

%==============================================================================
\section{Summary and Expected Impact}
%==============================================================================

Table~\ref{tab:summary} summarizes the implemented changes by category and expected effect.

\begin{table}[H]
\centering
\footnotesize
\renewcommand{\arraystretch}{1.15}
\begin{tabular}{@{}p{0.26\textwidth}p{0.17\textwidth}p{0.12\textwidth}p{0.35\textwidth}@{}}
\toprule
\textbf{Change} & \textbf{File(s)} & \textbf{Primary Axis} & \textbf{Expected Effect} \\
\midrule
Temporal guidance scale in sphere pipeline & \file{gen\_pano\_360.py}, \file{i2v\_sphere\_panorama\_pipeline.py} & Motion / consistency & Stronger temporal coherence and smoother motion. \\
Low-memory flag, reduced resolution and steps & \file{gen\_pano\_360.py} & Memory / throughput & Enables stable runs on 16\,GB GPUs with modest quality loss. \\
Latitude-adaptive theta sampling & \file{gen\_pano\_360.py} & Geometry / scene richness & Uniform angular coverage, fewer redundant polar views. \\
Cosine overlap and merge schedules & \file{gen\_pano\_360.py} & Temporal flicker & Fewer seam artifacts and sharper final details. \\
DDIM-aware re-noising & \file{scheduler.py} & Graphics / structure & Better structural preservation across noise-level changes. \\
Gaussian-smoothed latent upsampling & \file{diffusion\_utils.py} & Graphics & Less aliasing in upscaled latents. \\
Temporal Gaussian smoothing and batched save & \file{loop\_merge\_utils.py} & Temporal flicker / efficiency & Reduced flicker and faster decode/save. \\
Centered pixel grid and longitude normalization & \file{panorama\_tensor\_utils.py}, \file{ring\_panorama\_tensor\_utils.py} & Geometry & More accurate projections and better use of horizontal resolution. \\
Nearest-neighbor via rounding & Same as above & Geometry & Reduced reprojection bias. \\
Vectorized bilinear splatting with blend & Same as above & Geometry / code & Faster and more flexible multi-view fusion. \\
Adaptive tiled VAE encode with cosine weights & \file{i2v\_sphere\_panorama\_pipeline.py} & Image quality / memory & Smoother tile boundaries and scalable encoding to high resolutions. \\
Unified resize and normalization & \file{precast\_latent\_utils.py}, \file{tensor\_utils.py} & Image quality & Consistent conditioning and reduced artifacts. \\
Shift-window robustness in low-res mode & \file{i2v\_sphere\_panorama\_pipeline.py} & Stability & Prevents assertion failures in aggressive low-memory settings. \\
Q-Align / CLIP metric stubs & \file{quality\_metrics.py} & Evaluation & Infrastructure for future quality-aware selection and scoring. \\
\bottomrule
\end{tabular}
\caption{Summary of implemented changes and their intended impact.}
\label{tab:summary}
\end{table}

\section*{Conclusion}

The forked DynamicScaler codebase now implements a range of mathematically motivated and engineering-driven improvements:
geometric corrections, temporal guidance, smooth scheduling, adaptive tiling, and low-memory presets.
Together, these changes make the pipeline more robust on constrained GPUs, reduce visual artifacts (seams, flicker, aliasing), and provide hooks for future integration of learned quality metrics such as Q-Align and CLIP-based scores.

\end{document}

